\section{System model, scope and assumptions}\label{sec:model}

\subsection{Network model and scope}\label{sec:scope}
There are $n$ static sensor nodes $V=\{1,\dots,n\}$ and one sink $s$. The main setting is MADII's published setting~\cite{yang2025madii}: $n=100$ nodes uniform in a $500\times500$~m field, with the sink at the centre and initial energy $E_0=0.5$~J. Time is divided into rounds. In round $t$ node $i$ generates $k_i(t)$ packets, each forwarded to the sink along the round's routing tree $p_t:V\to V\cup\{s\}$. Relays forward every bit they receive, as in MADII (``no data compression''). \textbf{Scope.} The study concerns multi-hop tree routing without in-network aggregation. Protocols whose benefit comes mainly from aggregation at cluster heads (LEACH-family clustering) or that are specified for an aggregation setting (EBR-DA) are therefore not used as comparators. Comparing them with full forwarding would test them outside the setting they were designed for.

\subsection{Radio energy model}\label{sec:radio}
Transmitting $k$ bits over distance $d$ and receiving $k$ bits cost
\begin{align}
E_{\mathrm{TX}}(k,d) &= \begin{cases} kE_{\mathrm{elec}} + k\,\varepsilon_{\mathrm{fs}}d^2, & d<d_0,\\ kE_{\mathrm{elec}} + k\,\varepsilon_{\mathrm{amp}}d^4, & d\ge d_0,\end{cases}\label{eq:tx}\\
E_{\mathrm{RX}}(k) &= kE_{\mathrm{elec}},\qquad d_0=\sqrt{\varepsilon_{\mathrm{fs}}/\varepsilon_{\mathrm{amp}}}\approx 87.7~\mathrm{m},\label{eq:rx}
\end{align}
with $E_{\mathrm{elec}}=50$~nJ/bit, $\varepsilon_{\mathrm{fs}}=10$~pJ/bit/m$^2$ and $\varepsilon_{\mathrm{amp}}=0.0013$~pJ/bit/m$^4$. We write $e_{\mathrm{tx}}(i,j)=E_{\mathrm{TX}}(1,d_{ij})$ and $e_{\mathrm{rx}}=E_{\mathrm{elec}}$.

\subsection{Packet layer and control accounting}\label{sec:packet}
Every bit any router relies on is charged, and the same charges apply to every router (Table~\ref{tab:params}):
\begin{enumerate}[label=(\roman*), leftmargin=*]
  \item \textbf{Data packets} carry 4{,}000 payload bits plus a 16-bit residual-battery report from the source, so $L=4{,}016$ bits travel on every hop. Every energy-aware router, the LP planner and MADII's state upload need the batteries at the sink. The report is therefore charged to all routers, not assumed free.
  \item \textbf{The route advertisement} must encode every node's next hop. It has a 100-bit header plus $n\lceil\log_2(n+1)\rceil$ bits, so $L_c=800$ bits at $n=100$, and every alive node pays $E_{\mathrm{RX}}(L_c)$ per round. An earlier version of this study charged only 100 bits, which cannot encode 100 identifiers; Appendix~\ref{app:corrections} lists this and the other corrections.
  \item \textbf{Traffic rates} used by the routers are counted by the sink from the packets it receives, which costs no extra bits. The tiered-load variant (Section~\ref{sec:lest}) instead reports load in 2-bit tiers on an extra byte per packet plus a $2n$-bit snapshot in the advertisement. It is charged accordingly and is an ablation only.
\end{enumerate}
A relay that cannot afford a transmission dies within the round, and the packets routed through it in that round are lost. A source whose path failed retries directly to the sink if its battery allows. No MAC frames, acknowledgements or retransmissions are modelled in the main bench; Section~\ref{sec:lossy} adds lossy links, ARQ and contention as a probe.

\begin{table}[t]
\centering
\caption{Network, radio and packet parameters.}\label{tab:params}
\footnotesize
\begin{adjustbox}{max width=\linewidth}
\begin{tabular}{@{}L{0.46\linewidth}L{0.48\linewidth}@{}}
\toprule
Parameter & Value\\
\midrule
Nodes, field, sink & 100, uniform in $500\times500$~m, centre\\
Initial energy $E_0$ & 0.5~J (trace bench 1.0 and 1.5~J)\\
Data packet $L$ & 4{,}000 payload + 16 battery bits\\
Advertisement $L_c$ & $100 + n\lceil\log_2(n+1)\rceil$ (800 at $n=100$)\\
Radio constants & $E_{\mathrm{elec}}$ 50~nJ/bit; $\varepsilon_{\mathrm{fs}}$ 10~pJ/bit/m$^2$; $\varepsilon_{\mathrm{amp}}$ 0.0013~pJ/bit/m$^4$\\
Aggregation & none (full forwarding)\\
Tiered-load ablation & +8 bits per packet; +$2n$ bits per advertisement\\
\bottomrule
\end{tabular}
\end{adjustbox}
\end{table}

\subsection{Traffic}\label{sec:traffic}
A day is 24 rounds, and every router receives 3 days of pre-deployment history. In the daily-surge scenario (R), the 20\% of nodes nearest a random point send Poisson traffic with mean 4 packets per round during a 6-round evening and 0.5 otherwise, and all other nodes have mean 1. In the constant scenario (N), every node sends one packet per round, which is MADII's static setting. Bursts (B) adds to R, each round with probability 0.05, a burst in which a random node sends five times its rate for 4 rounds. Pattern shift (D) moves the hotspot to a new set of nodes after two days; it is a robustness scenario that nothing is tuned on. As a second, trace-driven traffic generator we convert Intel Berkeley Lab readings~\cite{intellab2004} into packets by send-on-delta reporting (one heartbeat per hour plus a packet per 0.5~$^\circ$C or 100~lux change) on the 54-mote layout scaled to the $500$~m field. This yields traffic with a different autocorrelation structure. It is not measured radio traffic.

\subsection{Assumptions}\label{sec:assumptions}
\begin{enumerate}[label=\textbf{A\arabic*}, leftmargin=*]
  \item Static nodes with identical initial energy; no harvesting or battery replacement.
  \item One static, mains-powered sink; its computation is not charged to sensors (its cost is reported separately in Section~\ref{sec:scaling}).
  \item No obstacles; link cost depends only on distance through~\eqref{eq:tx}--\eqref{eq:rx}.
  \item Any node can reach any other node and the sink by power control.
  \item Main bench: ideal symmetric links and collision-free scheduling. Section~\ref{sec:lossy} relaxes this.
  \item Only transmission and reception consume energy; sensing, processing and idle listening are not modelled.
  \item No aggregation or compression.
  \item Round-based operation: the sink computes each round's tree before the round and broadcasts it in the advertisement.
  \item The sink knows every node's position.
  \item The sink learns residual batteries from the 16-bit reports of the previous round (charged, item~(i)).
  \item The sink learns traffic only from packets it has received; only the oracle bound sees future traffic.
  \item Nodes fail only by depletion; a relay that cannot afford its transmission dies within the round.
  \item Every router sees the same deployments and traffic, so all comparisons are paired.
\end{enumerate}

\subsection{Lifetime, the bound and its tightness}\label{sec:problem}
Let $F$ be the round in which the first node dies; $T=F-1$ completed death-free rounds is the lifetime. For example, if the first death happens during round 152, then $F=152$ and $T=151$. Let $f_{ij}\ge0$ be the total bits sent on link $i\to j$. For constant traffic, the fractional-flow ceiling~\cite{chang2004maximum} is
\begin{align}
&\max_{T,\,f\ge0}\; T \quad\text{subject to, for every } i\in V,\nonumber\\
&\textstyle\sum_j f_{ij}-\sum_k f_{ki} = L\,T,\label{eq:lp-flow}\\
&\textstyle\sum_j f_{ij}\,e_{\mathrm{tx}}(i,j)+\sum_k f_{ki}\,e_{\mathrm{rx}} + T\,E_{\mathrm{RX}}(L_c) \le E_0.\label{eq:lp-energy}
\end{align}
For time-varying traffic the right-hand side of~\eqref{eq:lp-flow} becomes $L\sum_{t<T}k_i(t)$, and $T^*$ is the largest $T$ for which the LP is feasible, found by bisection (an oracle that knows the future). Any sequence of routing trees that keeps all nodes alive for $T$ rounds gives a feasible flow, so $T\le T^*$. We report every router as $100\,T/T^*$. This is a percentage of a \emph{relaxation}, not of the best achievable sequence of trees. Its looseness is bounded empirically in Section~\ref{sec:tight}: on every case, the best realised router certifies that the achievable optimum lies between that router's lifetime and $T^*$.
